%Data institusi

%\input{Dozentdaten}

%\renewcommand{\fachbereich}{Fachbereich}

%\renewcommand{\dozent}{Prof. Dr. . }

%Data ujian

\renewcommand{\klausurgebiet}{ }

\renewcommand{\klausurtyp}{ }

\renewcommand{\klausurdatum}{ . 20}

\klausurvorspann {\fachbereich} {\klausurdatum} {\dozent} {\klausurgebiet} {\klausurtyp}

%Data untuk tabel nilai berikut

\renewcommand{\aeins}{  3  }

\renewcommand{\azwei}{  3   }

\renewcommand{\adrei}{  8  }

\renewcommand{\avier}{  7  }

\renewcommand{\afuenf}{  8  }

\renewcommand{\asechs}{  3  }

\renewcommand{\asieben}{  5  }

\renewcommand{\aacht}{  5  }

\renewcommand{\aneun}{  6  }

\renewcommand{\azehn}{  4  }

\renewcommand{\aelf}{  12  }

\renewcommand{\azwoelf}{  64  }

\renewcommand{\adreizehn}{     }

\renewcommand{\avierzehn}{    }

\renewcommand{\afuenfzehn}{    }

\renewcommand{\asechzehn}{    }

\renewcommand{\asiebzehn}{    }

\renewcommand{\aachtzehn}{    }

\renewcommand{\aneunzehn}{    }

\renewcommand{\azwanzig}{    }

\renewcommand{\aeinundzwanzig}{    }

\renewcommand{\azweiundzwanzig}{    }

\renewcommand{\adreiundzwanzig}{    }

\renewcommand{\avierundzwanzig}{    }

\renewcommand{\afuenfundzwanzig}{    }

\renewcommand{\asechsundzwanzig}{    }

\punktetabelleelf

\klausurnote

\newpage

\setcounter{bagian}{0}



\inputaufgabeklausurloesung
{3}

{

Definisikan istilah-istilah berikut
\zusatzklammer {yang dicetak miring} {} {.}
\aufzaehlungsechs{Suatu
\stichwort {medan normal} {}
$N$ pada hiperpermukaan diferensiabel

\mavergleichskette
{\vergleichskette
{ Y
}
{ \subseteq }{ \R^n
}
{  }{
}
{    }{
}
{   }{
}
}
{}{}{.}

}{Suatu ruang topologis yang
\stichwort {parakompak}{}
.

}{
\stichwort {Ruang kotangen} {}
di suatu titik

\mavergleichskette
{\vergleichskette
{  P
}
{ \in }{ M
}
{  }{}
{    }{}
{   }{}
}
{}{}{}
dari
\definitionsverweis {manifold diferensiabel}{}{}
$M$.

}{Suatu
\stichwort {titik reguler} {}

\mathl{Q \in L}{} bagi
\definitionsverweis {pemetaan diferensiabel}{}{}
\maabbdisp {\varphi} { L } { M
} {}
antara
\definitionsverweis {manifold-manifold diferensiabel}{}{}
\mathkor {} {L} {dan} {M} {.}

}{
\stichwort {Produk} {}
dari manifold-manifold diferensiabel
\mathkor {} {M} {dan} {N} {.}

}{Suatu
\stichwort {koneksi} {}
pada bundel vektor diferensiabel
\maabb {} {E} {M
} {}
di atas manifold diferensiabel $M$.
}

}
{

\aufzaehlungsechs{Suatu
medan normal
pada $Y$ adalah objek yang didefinisikan pada suatu lingkungan terbuka

\mavergleichskette
{\vergleichskette
{ Y
}
{ \subseteq }{ U
}
{ \subseteq }{ \R^n
}
{ }{
}
{ }{
}
}
{}{}{}
dan bersifat
\definitionsverweis {kontinu}{}{}
sebagai suatu
\definitionsverweis {medan vektor}{}{}
\maabb {F} {U} {\R^n
} {}
yang memenuhi

\mavergleichskettedisp
{\vergleichskette
{ F(P)
}
{ \in} { N_PY
}
{ } {
}
{ } {
}
{ } {
}
}
{}{}{}
untuk setiap

\mavergleichskette
{\vergleichskette
{ P
}
{ \in }{ Y
}
{  }{
}
{    }{
}
{   }{
}
}
{}{}{.}
}{Suatu
\definitionsverweis {ruang topologis}{}{}
$X$ disebut kompak terhadap penutup jika untuk setiap penutup terbuka

\mathdisp {X= \bigcup_{i \in I} U_i \, \, \, \text{ dengan } U_i \text{ terbuka dan himpunan indeks sembarang }} {  }

terdapat suatu himpunan bagian berhingga
\mathl{J \subseteq I}{} sedemikian sehingga

\mavergleichskettedisp
{\vergleichskette
{ X
}
{ =} {\bigcup_{i \in J} U_i
}
{ } {
}
{ } {
}
{ } {
}
}
{}{}{}
berlaku.
}{Ruang kotangen di $P$ adalah
\definitionsverweis {ruang dual}{}{}
dari
\definitionsverweis {ruang singgung}{}{}

\mathl{T_PM}{} di $P$.
}{Titik

\mavergleichskette
{\vergleichskette
{ Q
}
{ \in }{ L
}
{ }{
}
{ }{
}
{ }{
}
}
{}{}{}
disebut reguler bagi $\varphi$ jika
\definitionsverweis {pemetaan tangen}{}{}
\maabbdisp {T_Q(\varphi)} {T_QL} {T_{\varphi(Q)}M
} {}
pada titik $Q$ mempunyai
\definitionsverweis {rank maksimal}{}{.}
}{Misalkan
\mathkor {} {(U_i,U_i',\alpha_i, i \in I)} {dan} {(V_j,V_j',\beta_j, j \in J)} {}
merupakan
\definitionsverweis {atlas-atlas}{}{}
dari
\mathkor {} {M} {dan} {N} {.}
Maka
\definitionsverweis {ruang produk}{}{}

\mathl{M \times N}{} yang dilengkapi dengan
\definitionsverweis {peta-peta}{}{}
\maabbdisp {\alpha_i \times \beta_j} {U_i \times V_j } { U_i' \times V_j'
} {}
\zusatzklammer {dengan \mathlk{(i,j) \in I \times J}{} dan \mathlk{U_i' \times V_j' \subseteq \R^m \times \R^n}{}} {} {}
disebut produk manifold
\mathkor {} {M} {dan} {N} {.}
}{Suatu
koneksi
pada $E$ adalah dekomposisi jumlah langsung dari
\definitionsverweis {bundel tangen}{}{}

\mavergleichskette
{\vergleichskette
{ TE
}
{ = }{ V \oplus H
}
{ }{
}
{ }{
}
{ }{
}
}
{}{}{}
menjadi dua subbundel vektor
\mathkor {} {V} {dan} {H} {,}
dengan

\mavergleichskette
{\vergleichskette
{ V
}
{ \subseteq }{ TE
}
{ }{
}
{ }{
}
{ }{
}
}
{}{}{}
merupakan
\definitionsverweis {bundel vertikal}{}{.}
}

 }



\inputaufgabeklausurloesung
{3}

{

Nyatakan teorema-teorema berikut.
\aufzaehlungdrei{
\stichwort {Teorema tentang hubungan antara kelengkungan dan vektor normal satuan suatu kurva bidang berparameter panjang busur} {.}}{
\stichwort {Theorema egregium} {.}}{
\stichwort {Teorema Green} {}
untuk luas.}

}
{

\aufzaehlungdrei{\faktsituation {Misalkan
\maabb {\gamma} { I } {\R^2
} {}
suatu
\definitionsverweis {kurva}{}{}
\definitionsverweis {berparameter panjang busur}{}{}
yang dua kali
\definitionsverweis {terdiferensialkan secara kontinu}{}{.}}
\faktfolgerung {Maka
\definitionsverweis {kelengkungan}{}{}
dari $\gamma$ di $t$ adalah

\mavergleichskettedisp
{\vergleichskette
{ \kappa(t)
}
{ =} { \left\langle \gamma^{\prime \prime}(t) , N(t) \right\rangle
}
{ = } { \pm \Vert { \gamma^{\prime \prime}(t) } \Vert
}
{ } {
}
{ } {
}
}
{}{}{,}
dengan

\mavergleichskettedisp
{\vergleichskette
{ N(t)
}
{ =} { \begin{pmatrix} -\gamma_2'(t) \\ \gamma_1'(t) \end{pmatrix}
}
{ } {
}
{ } {
}
{ } {
}
}
{}{}{}
suatu vektor normal satuan di $\gamma(t)$.}
\faktzusatz {}
\faktzusatz {}}{\faktsituation {Misalkan

\mavergleichskette
{\vergleichskette
{ Y
}
{ \subseteq }{ \R^3
}
{  }{
}
{    }{
}
{   }{
}
}
{}{}{}
suatu
\definitionsverweis {permukaan berorientasi}{}{}
dan misalkan
\maabbdisp {\varphi} { V } { Y
} {,}

\mavergleichskette
{\vergleichskette
{ V
}
{ \subseteq }{ \R^2
}
{  }{
}
{    }{
}
{   }{
}
}
{}{}{,}
suatu parametrisasi lokal $Y$ yang dua kali terdiferensialkan, dengan parameter $u,v$. Misalkan
\mathl{\left(  g_{ij}  \right)}{}
\definitionsverweis {matriks fundamental pertama}{}{}
pada $V$.}
\faktfolgerung {Maka, dengan menggunakan
\definitionsverweis {simbol-simbol Christoffel}{}{,}
\definitionsverweis {kelengkungan Gauss}{}{}
memenuhi hubungan

\mavergleichskettedisphandlinks
{\vergleichskettedisphandlinks
{ K
}
{ =} { { \frac{ 1 }{ g_{11} } } \left( \partial_2 \Gamma_{11}^2 - \partial_1 \Gamma_{12}^2 + \Gamma^1_{11} \Gamma^2_{12} + \Gamma_{11}^2 \Gamma_{22}^2 - \Gamma_{12}^1 \Gamma_{11}^2- \left( \Gamma_{12}^2 \right)^2 \right)
}
{ } {
}
{ } {
}
{ } {
}
}
{}{}{.}}
\faktzusatz {}
\faktzusatz {}}{Misalkan
\mathl{M \subseteq \R^2}{} suatu
\definitionsverweis {manifold dengan batas}{}{}
yang
\definitionsverweis {kompak}{}{}

\mathl{\delta M}{.} Maka luas $M$ adalah

\mavergleichskettedisp
{\vergleichskette
{  \lambda^2(M)
}
{ =} {
\int_{ \partial M  }  x dy
}
{ =} { -
\int_{ \partial M  }  y dx
}
{ } {
}
{ } {
}
}
{}{}{.}}

 }



\inputaufgabeklausurloesung
{8}

{

Misalkan

\mavergleichskettedisp
{\vergleichskette
{ f(x,y)
}
{ =} { ax^2+by^2+cxy
}
{ } {
}
{ } {
}
{ } {
}
}
{}{}{.}
Tentukan
\definitionsverweis {kelengkungan-kelengkungan utama}{}{}
dan
\definitionsverweis {arah-arah kelengkungan utama}{}{}
dari
\definitionsverweis {grafik}{}{}
$f$ di titik asal.

}
{

Kita menggunakan
Lemma 4.9 (Differentialgeometrie (Osnabrück 2023)),
dengan faktor pendahulu di titik nol sama dengan $1$, sehingga
\definitionsverweis {pemetaan Weingarten}{}{}
dideskripsikan langsung oleh
\definitionsverweis {matriks Hesse}{}{}
. Matriks Hesse fungsi tersebut adalah

\mathdisp {\begin{pmatrix} 2a & c \\ c & 2b \end{pmatrix}} { , }

dan
\definitionsverweis {polinomial karakteristik}{}{}
matriks tersebut adalah

\mavergleichskettealign
{\vergleichskettealign
{ (t-2a)(t-2b)-c^2
}
{ =} { t^2 -2(a+b)t +4ab -c^2
}
{ =} { \left(  t-(a+b)  \right)^2 - (a+b)^2 +4ab -c^2
}
{ =} { \left(  t-(a+b)  \right)^2  - a^2 -b^2 +2ab -c^2
}
{ =} { \left(  t-(a+b)  \right)^2  - (a-b)^2 -c^2
}
}
{}
{}{.}
Dengan demikian, nilai-nilai eigen
\zusatzklammer {yang merupakan kelengkungan-kelengkungan utama} {} {}
adalah

\mavergleichskettedisp
{\vergleichskette
{ \lambda_{1,2}
}
{ =} { \pm \sqrt{ (a-b)^2+c^2 } +a+b
}
{ } {
}
{ } {
}
{ } {
}
}
{}{}{.}
Dalam kasus

\mavergleichskette
{\vergleichskette
{ a
}
{ = }{ b
}
{  }{
}
{    }{
}
{   }{
}
}
{}{}{}
dan

\mavergleichskette
{\vergleichskette
{c
}
{ = }{ 0
}
{ }{
}
{ }{
}
{ }{
}
}
{}{}{}
pemetaan Weingarten adalah perkalian dengan $2a$ dan setiap vektor merupakan vektor eigen; mulai sekarang kasus ini dikecualikan. Untuk

\mavergleichskettedisp
{\vergleichskette
{ \lambda_{1}
}
{ =} { \sqrt{ (a-b)^2+c^2 } +a+b
}
{ } {
}
{ } {
}
{ } {
}
}
{}{}{}
perlu ditentukan kernel matriks

\mathdisp {\begin{pmatrix}  \sqrt{(a-b)^2 +c^2} + b-a  &   -c  \\  -c &  \sqrt{(a-b)^2 +c^2} + a-b \end{pmatrix}} {  }

dan kernel tersebut adalah

\mathdisp {\R \begin{pmatrix}  \sqrt{(a-b)^2 +c^2} + a-b  \\c   \end{pmatrix}} { . }

Jadi, salah satu arah kelengkungan utama adalah
\mathl{\begin{pmatrix}  \sqrt{(a-b)^2 +c^2} + a-b  \\c   \end{pmatrix}}{} dengan kelengkungan utama
\mathl{\sqrt{ (a-b)^2+c^2 } +a+b}{.}

Dengan cara yang sama diperoleh arah kelengkungan utama
\mathl{\begin{pmatrix} -c \\ \sqrt{(a-b)^2 +c^2} + a-b \end{pmatrix}}{} dengan kelengkungan utama
\mathl{-\sqrt{ (a-b)^2+c^2 } +a+b}{.}

 }



\inputaufgabeklausurloesung
{7}

{

Buktikan teorema eksistensi untuk transpor paralel pada suatu hiperpermukaan

\mavergleichskette
{\vergleichskette
{ Y
}
{ \subseteq }{ \R^n
}
{  }{
}
{    }{
}
{   }{
}
}
{}{}{.}

}
{

Kita mencari suatu pemetaan diferensiabel
\maabbdisp {F} { I } { \R^n
} {,}
yang memenuhi syarat-syarat
\aufzaehlungdrei{
\mavergleichskettedisp
{\vergleichskette
{ F(t)
}
{ \in} { T_{\gamma(t)} Y
}
{ } {
}
{ } {
}
{ } {
}
}
{}{}{}
untuk setiap

\mavergleichskette
{\vergleichskette
{ t
}
{ \in }{ I
}
{  }{
}
{    }{
}
{   }{
}
}
{}{}{,}
}{
\mavergleichskettedisp
{\vergleichskette
{ F'(t)
}
{ \in} { N_{\gamma(t)} Y
}
{ } {
}
{ } {
}
{ } {
}
}
{}{}{}
untuk setiap

\mavergleichskette
{\vergleichskette
{ t
}
{ \in }{ I
}
{ }{
}
{ }{
}
{ }{
}
}
{}{}{,}
}{
\mavergleichskettedisp
{\vergleichskette
{ F (t_0)
}
{ =} { v
}
{ } {
}
{ } {
}
{ } {
}
}
{}{}{}
}
Syarat pertama berarti bahwa $F(t)$ tegak lurus terhadap vektor normal satuan $N (\gamma(t))$, sehingga

\mavergleichskettedisp
{\vergleichskette
{ \left\langle F( t) , N( \gamma(t)) \right\rangle
}
{ =} { 0
}
{ } {
}
{ } {
}
{ } {
}
}
{}{}{.}
Oleh karena itu, juga berlaku

\mavergleichskettedisp
{\vergleichskette
{ 0
}
{ =} { \left\langle F(t) , N( \gamma(t)) \right\rangle'
}
{ =} { \left\langle F'(t) , N( \gamma(t)) \right\rangle + \left\langle F(t) , (N( \gamma(t)))' \right\rangle
}
{ } {
}
{ } {
}
}
{}{}{.}
Syarat kedua, yakni bahwa $F'(t)$ terletak di ruang normal, berarti

\mavergleichskettedisp
{\vergleichskette
{ F'(t) - \left\langle F'( t) , N( \gamma(t)) \right\rangle N(\gamma(t) )
}
{ =} { 0
}
{ } {
}
{ } {
}
{ } {
}
}
{}{}{}
untuk setiap

\mavergleichskette
{\vergleichskette
{ t
}
{ \in }{ I
}
{  }{
}
{    }{
}
{   }{
}
}
{}{}{,}
Dengan menggunakan persamaan sebelumnya, syarat ini dapat diubah menjadi

\mavergleichskettedisp
{\vergleichskette
{ F'(t) + \left\langle F( t) , (N( \gamma(t)))' \right\rangle N(\gamma(t) )
}
{ =} { 0
}
{ } {
}
{ } {
}
{ } {
}
}
{}{}{}
untuk setiap

\mavergleichskette
{\vergleichskette
{ t
}
{ \in }{ I
}
{  }{
}
{    }{
}
{   }{
}
}
{}{}{}
atau menjadi

\mavergleichskettedisp
{\vergleichskette
{ F'(t)
}
{ =} { - \left\langle F( t) ,  (N( \gamma(t)))'   \right\rangle N(\gamma(t) )
}
{ } {
}
{ } {
}
{ } {
}
}
{}{}{}
untuk setiap

\mavergleichskette
{\vergleichskette
{ t
}
{ \in }{ I
}
{  }{
}
{    }{
}
{   }{
}
}
{}{}{}
. Ini adalah persamaan diferensial linear biasa orde pertama untuk $F$. Bersama syarat awal

\mavergleichskette
{\vergleichskette
{ F(t_0)
}
{ = }{ v
}
{ }{
}
{ }{
}
{ }{
}
}
{}{}{}
menurut
Satz 56.2 (Analysis (Osnabrück 2021-2023))
persamaan tersebut mempunyai solusi yang tunggal. Jadi, sistem syarat semula mempunyai paling banyak satu solusi. Jika $F$ adalah solusi tunggal persamaan diferensial itu, maka pemetaan tersebut memang merupakan solusi sistem syarat semula.

 }



\inputaufgabeklausurloesung
{8}

{

Tunjukkan bahwa suatu
\definitionsverweis {manifold diferensiabel}{}{}
$M$ berdimensi

\mavergleichskette
{\vergleichskette
{ n
}
{ \geq }{ 1
}
{  }{
}
{    }{
}
{   }{
}
}
{}{}{}
mempunyai tak hingga banyak
\definitionsverweis {difeomorfisme}{}{}
\maabbdisp {\varphi} { M } { M
} {}
.

}
{

Kita meninjau fungsi-fungsi berbentuk
\maabbeledisp {\psi_c} {[0,1]} {[0,1]
} { x } { x + \varphi_c (x)
} {,}
dengan

\mavergleichskettedisp
{\vergleichskette
{  \varphi_c(x)
}
{ =} { \begin{cases}  0 \text{ jika } x \leq { \frac{   1  }{  3 } } \, , \\  c \left(  x- { \frac{   1  }{  3 } }   \right)^2 \left(  x- { \frac{   2 }{  3 } }   \right)^2 \text{ jika } { \frac{   1  }{  3 } } < x < { \frac{   2 }{  3 } }  \, ,  \\  0 \text{ jika } x \geq { \frac{   2 }{  3 } }  \, .  \end{cases}
}
{ } {
}
{ } {
}
{ } {
}
}
{}{}{}
dengan parameter

\mavergleichskette
{\vergleichskette
{ c
}
{ \in }{  \R_{\geq 0}
}
{  }{
}
{    }{
}
{   }{
}
}
{}{}{.}
Fungsi $\varphi_c$ kontinu dan juga terdiferensialkan secara kontinu karena akar-akar polinomial tersebut merupakan akar ganda. Dengan demikian, $\psi_c$ juga terdiferensialkan secara kontinu. Untuk $c$ yang cukup kecil, $\psi_c$ meningkat secara ketat dan mendefinisikan suatu pemetaan bijektif yang terdiferensialkan secara kontinu dari interval satuan ke dirinya sendiri, serta merupakan identitas pada sepertiga bawah dan sepertiga atas.

Dengan $\psi_c$, kita mendefinisikan difeomorfisme bola satuan terbuka ke dirinya sendiri melalui
\maabbeledisp {} { U \left(   0 , 1  \right) } { U \left(   0 , 1  \right)
} { \begin{pmatrix}  x_1  \\\vdots \\  x_n   \end{pmatrix} } { \psi_c \left(  \sqrt{  x_1^2 + \cdots + x_n^2 }  \right) \begin{pmatrix}  x_1  \\\vdots \\  x_n   \end{pmatrix}
} {.}
Jadi, titik-titik diregangkan bergantung pada jari-jarinya, sedangkan pemetaan tersebut merupakan identitas pada
\mathl{U \left(   0 , { \frac{   1  }{  3 } }   \right)}{} dan pada
\mathl{U \left(   0 , 1  \right) \setminus  U \left(   0 ,   { \frac{   2 }{  3 } }   \right)}{}.

Jika sekarang $M$ suatu manifold berdimensi $\geq 1$, kita memilih suatu peta
\maabbdisp {\alpha} { U } {V
} {}
dengan

\mavergleichskette
{\vergleichskette
{ V
}
{ \subseteq }{ \R^n
}
{ }{
}
{ }{
}
{ }{
}
}
{}{}{,}
dan, dengan
\zusatzklammer {memperkecil daerahnya} {} {,}
kita dapat menganggap bahwa $V$
\zusatzklammer {merupakan suatu bola terbuka dan kemudian} {} {}
adalah bola satuan terbuka. Difeomorfisme yang dikonstruksi pada $V$ menghasilkan difeomorfisme pada $U$. Karena difeomorfisme ini merupakan identitas pada bagian luar bola
\zusatzklammer {mulai dari jari-jari ${ \frac{   2 }{  3 } }$} {} {,}
kita dapat memperluasnya secara diferensiabel ke $M$ dengan menggunakan identitas pada $M \setminus \alpha^{-1} \left(  U \left(   0 , { \frac{   2 }{  3 } }   \right)  \right)$.

 }



\inputaufgabeklausurloesung
{3}

{

Misalkan $M$ suatu
\definitionsverweis {manifold diferensiabel}{}{}
dan
\maabbdisp {f} { M } {\R
} {}
suatu
\definitionsverweis {fungsi terdiferensialkan secara kontinu}{}{.}
Misalkan di titik

\mavergleichskette
{\vergleichskette
{ P
}
{ \in }{ M
}
{  }{
}
{    }{
}
{   }{
}
}
{}{}{}
fungsi tersebut mempunyai
\definitionsverweis {ekstremum lokal}{}{.}
Tunjukkan bahwa

\mavergleichskettedisp
{\vergleichskette
{ (df)_P
}
{ =} { 0
}
{ } {
}
{ } {
}
{ } {
}
}
{}{}{.}

}
{

Misalkan

\mavergleichskette
{\vergleichskette
{ P
}
{ \in }{ U
}
{ \subseteq }{ M
}
{ }{
}
{ }{
}
}
{}{}{}
suatu
\definitionsverweis {daerah peta}{}{}
dengan peta
\maabb {\alpha} { U } {V \subseteq \R^n
} {.}
Maka
\mathl{f \circ \alpha^{-1}}{} juga mempunyai ekstremum lokal di $\alpha(P)$.
Oleh karena itu, menurut
Fakt ***** (2)
\definitionsverweis {diferensial total}{}{}
\maabbdisp {\left(D f \circ \alpha^{-1} \right)_{\alpha(P) }} {\R^n } {\R
} {}
adalah pemetaan nol. Karena itu,

\mavergleichskette
{\vergleichskette
{ (df)_P
}
{ = }{ 0
}
{ }{
}
{ }{
}
{ }{
}
}
{}{}{.}

 }



\inputaufgabeklausurloesung
{5}

{

Kita tinjau
$1$-\definitionsverweis {bentuk diferensial}{}{}

\mavergleichskettedisp
{\vergleichskette
{ \omega
}
{ =} { -vdu+udv
}
{ } {
}
{ } {
}
{ } {
}
}
{}{}{}
pada
$1$-\definitionsverweis {bola}{}{}

\mavergleichskette
{\vergleichskette
{ S^1
}
{ \subset }{ \R^2
}
{  }{
}
{    }{
}
{   }{
}
}
{}{}{,}
dengan koordinat ruang sekitar $\R^2$ dinotasikan oleh
\mathkor {} {u} {dan} {v} {.}
Tentukan
\mathl{\pi^* \omega}{} di bawah
\definitionsverweis {pemetaan terdiferensialkan secara kontinu}{}{}
\maabbeledisp {\pi} {\R^2 \setminus \{0\}} { S^{1}
} {(x , y) } { { \frac{   1  }{  \sqrt{x^2 + y^2}  } } (x , y)
} {.}

}
{

Berlaku

\mavergleichskettealign
{\vergleichskettealign
{ \pi^*du
}
{ =} { d (u \circ \pi)
}
{ =} { d   \left(  { \frac{   x  }{  \sqrt{x^2 + y^2}  } }  \right)
}
{ =} { { \frac{   \partial   \left(  { \frac{   x  }{  \sqrt{x^2 + y^2}  } }  \right)    }{  \partial   x  } }   dx + { \frac{   \partial   \left(  { \frac{   x  }{  \sqrt{x^2 + y^2}  } }  \right)    }{  \partial   y  } } dy
}
{ =} { { \frac{   \sqrt{x^2+y^2}     -x^2 (x^2 + y^2)^{-1/2}  }{  x^2+y^2  } }  dx   - { \frac{   xy }{  \sqrt{x^2+y^2}^3  } }  dy
}
}
{
\vergleichskettefortsetzungalign
{ =} { { \frac{   x^2+y^2 -x^2    }{  \sqrt{ x^2+y^2}^3  } }  dx  - { \frac{   xy }{  \sqrt{x^2+y^2}^3  } }  dy
}
{ =} { { \frac{   y^2  }{  \sqrt{ x^2+y^2}^3  } }  dx  - { \frac{   xy }{  \sqrt{x^2+y^2}^3  } }  dy
}
{ } {}
{ } {}
}
{}{}
dan

\mavergleichskettealign
{\vergleichskettealign
{ \pi^*dv
}
{ =} { d (v \circ \pi)
}
{ =} { d \left(  { \frac{   y  }{  \sqrt{x^2 + y^2}  } }  \right)
}
{ =} { { \frac{   \partial   \left(  { \frac{   y  }{  \sqrt{x^2 + y^2}  } }  \right)   }{  \partial   x  } } dx + { \frac{   \partial   \left(  { \frac{   y  }{  \sqrt{x^2 + y^2}  } }  \right)    }{  \partial   y  } } dy
}
{ =} { { \frac{   -xy }{  \sqrt{x^2+y^2}^3  } } dx + { \frac{   \sqrt{x^2+y^2} -y^2 (x^2 + y^2)^{-1/2}  }{  x^2+y^2  } } dy
}
}
{
\vergleichskettefortsetzungalign
{ =} { { \frac{   -xy }{  \sqrt{x^2+y^2}^3  } } dx + { \frac{   x^2  }{  \sqrt{x^2+y^2}^3  } }  dy
}
{ } {
}
{ } {}
{ } {}
}
{}{.}
Dengan demikian,

\mavergleichskettealigndrucklinks
{\vergleichskettealigndrucklinks
{ \pi^*(-vdu+udv)
}
{ =} { - (v \circ \pi)  \pi^*du + (u \circ \pi) \pi^*dv
}
{ =} { - { \frac{   y  }{  \sqrt{x^2+y^2}  } } \left(  { \frac{   y^2  }{  \sqrt{ x^2+y^2}^3  } }  dx  - { \frac{   xy }{  \sqrt{x^2+y^2}^3  } }  dy   \right) + { \frac{   x  }{  \sqrt{x^2+y^2}  } }  \left(  { \frac{   -xy }{  \sqrt{x^2+y^2}^3  } } dx + { \frac{   x^2  }{  \sqrt{x^2+y^2}^3  } } dy  \right)
}
{ =} { { \frac{   1 }{  \left(  x^2+y^2  \right)^2 } }  \left(  \left(  -y^3-x^2y  \right) dx + \left(  xy^2+x^3  \right) dy \right)
}
{ =} { { \frac{   1 }{  x^2+y^2  } }  \left(  -  y dx + x dy \right)
}
}
{}{}{.}

 }



\inputaufgabeklausurloesung
{5}

{

Buktikan teorema tentang perhitungan volume kanonik pada manifold Riemann $M$.

}
{

Menurut
Definition  .
kita harus menghitung bentuk diferensial
\mathl{\left(  \alpha^{-1}  \right)^*\omega}{} untuk setiap titik

\mavergleichskette
{\vergleichskette
{ Q
}
{ \in }{ V
}
{  }{
}
{    }{
}
{   }{
}
}
{}{}{}
. Bentuk ini mempunyai bentuk
\mathl{c_Q dx_1 \wedge \ldots \wedge dx_n}{} dan ditentukan oleh nilainya pada
\mathl{e_1 \wedge \ldots \wedge e_n}{}. Berlaku

\mavergleichskettedisphandlinks
{\vergleichskettedisphandlinks
{ \left(  \alpha^{-1}  \right)^*\omega \left(  Q,  e_1 \wedge \ldots \wedge e_n  \right)
}
{ =} { \omega \left(  \alpha^{-1} (Q) ,T_Q \left(  \alpha^{-1}  \right)(e_1) \wedge \ldots \wedge T_Q \left(  \alpha^{-1}  \right)( e_n)  \right)
}
{ } {
}
{ } {
}
{ } {
}
}
{}{}{.}
Menurut definisi matriks fundamental metrik,

\mavergleichskettedisp
{\vergleichskette
{ g_{ij}(Q)
}
{ =} { \left\langle T_Q \left( \alpha^{-1} \right) (e_i) , T_Q \left( \alpha^{-1} \right) (e_j) \right\rangle_{\alpha^{-1}(Q)}
}
{ } {
}
{ } {
}
{ } {
}
}
{}{}{.}
Menurut
Satz 7.8 (Maß- und Integrationstheorie (Osnabrück 2022-2023)),
berlaku

\mavergleichskettealigndrucklinks
{\vergleichskettealigndrucklinks
{ \omega \left(  \alpha^{-1}(Q), T_Q \left(  \alpha^{-1}  \right) (e_1) \wedge \ldots \wedge T_Q \left(  \alpha^{-1}  \right) ( e_n)  \right)
}
{ =} { \left(  \det  \left(  \left\langle T_Q \left(  \alpha^{-1}  \right)  (e_i)  ,  T_Q \left(  \alpha^{-1}  \right) (e_j)   \right\rangle  \right)_{1 \leq i,j \leq n}  \right)^{1/2}
}
{ =} { \left(  \det  \left(  g_{ij}(Q)   \right)_{1 \leq i,j \leq n}  \right)^{1/2}
}
{ =} { \sqrt{g(Q)}
}
{ } {
}
}
{}{}{.}

 }



\inputaufgabeklausurloesung
{6 (1+1+3+1)}

{

Jelaskan mengapa $\R^2$ tidak membawa struktur
\definitionsverweis {manifold berbatas}{}{}
sedemikian sehingga subhimpunan $T$ yang diberikan merupakan
\definitionsverweis {batas}{}{.}
\aufzaehlungvierabc{
\mavergleichskette
{\vergleichskette
{ T
}
{ = }{ {]0,1[}
}
{  }{
}
{    }{
}
{   }{
}
}
{}{}{,}
dengan interval tersebut terletak pada sumbu $x$.
}{
\mavergleichskette
{\vergleichskette
{ T
}
{ = }{ [0,1]
}
{  }{
}
{    }{
}
{   }{
}
}
{}{}{,}
dengan interval tersebut terletak pada sumbu $x$.
}{
\mavergleichskette
{\vergleichskette
{ T
}
{ = }{ \R
}
{  }{
}
{    }{
}
{   }{
}
}
{}{}{,}
dengan $\R$ adalah sumbu $x$.
}{
\mavergleichskette
{\vergleichskette
{ T
}
{ = }{ S^1
}
{  }{
}
{    }{
}
{   }{
}
}
{}{}{.}
}

}
{

a) Batas
\mathl{\partial M}{} dari suatu manifold dengan batas, menurut
Lemma 21.3 (Differentialgeometrie (Osnabrück 2023))  (2)
\definitionsverweis {tertutup}{}{,}
sedangkan interval terbuka sebagai himpunan bagian dari $\R^2$ tidak tertutup.

b) Batas
\mathl{\partial M}{} dari suatu manifold dengan batas, menurut
Satz 21.4 (Differentialgeometrie (Osnabrück 2023)),
merupakan manifold tanpa batas, sedangkan interval tertutup mempunyai titik-titik batas.

c) Andaikan $\R^2$ merupakan suatu manifold dengan batas
\mathl{\R = \R \times \{0\}}{}. Untuk
\mathl{P \in \R}{}, terdapat suatu lingkungan terbuka
\mathl{P\in U}{} dan suatu homeomorfisme
\maabbdisp {\alpha} { U } {V
} {}
dengan suatu himpunan terbuka $V \subseteq H$, dengan $H$ menyatakan suatu setengah bidang. Kita dapat mengganti $V$ dengan lingkungan setengah bola terbuka yang lebih kecil di sekitar $P$. Menurut
Lemma 21.3 (Differentialgeometrie (Osnabrück 2023))  (2)
peta seperti ini memetakan titik batas ke titik batas; artinya,

\mavergleichskettedisp
{\vergleichskette
{\alpha ( U \cap \R )
}
{ = } { V \cap \delta H
}
{ } {
}
{ } {
}
{ } {
}
}
{}{}{.}
Dengan demikian, diperoleh pula suatu homeomorfisme di antara komplemennya, yaitu antara
\mathkor {} {U \setminus U \cap \R} {dan} {V \setminus V \cap \delta H} {.}
Lingkungan setengah bola di sebelah kanan
\definitionsverweis {terhubung lintasan}{}{,}
sedangkan himpunan di sebelah kiri merupakan gabungan saling lepas dari irisannya dengan setengah bagian positif dan negatif. Kedua irisan tersebut terbuka dan tidak kosong karena $U$ memuat suatu lingkungan bola dari $P$. Jadi, himpunan di sebelah kiri tidak terhubung dan homeomorfisme semacam itu tidak mungkin ada.

d) Sama seperti c).

 }



\inputaufgabeklausurloesung
{4}

{

Misalkan $\omega$ suatu
$n-1$-\definitionsverweis {bentuk diferensial}{}{}
yang terdiferensialkan secara kontinu pada
\definitionsverweis {himpunan terbuka}{}{}

\mavergleichskette
{\vergleichskette
{ U
}
{ \subseteq }{ \R^n
}
{  }{
}
{    }{
}
{   }{
}
}
{}{}{}
dan misalkan
\definitionsverweis {turunan eksterior}{}{}
sama dengan

\mavergleichskettedisp
{\vergleichskette
{d \omega
}
{ =} { fdx_1 \wedge \ldots \wedge dx_n
}
{ } {
}
{ } {
}
{ } {
}
}
{}{}{}
dengan suatu fungsi
\maabb {f} { U } {\R
} {.}
Tunjukkan untuk

\mavergleichskette
{\vergleichskette
{ P
}
{ \in }{ U
}
{  }{
}
{    }{
}
{   }{
}
}
{}{}{}
kesamaan

\mavergleichskettedisp
{\vergleichskette
{ f(P)
}
{ =} { { \frac{   1  }{  \beta_n  } } \cdot \operatorname{lim}_{   \epsilon \rightarrow  0  } \, { \frac{   1  }{  \epsilon^n } } \int_{S^{n-1}(P, \epsilon)} \omega
}
{ } {
}
{ } {
}
{ } {
}
}
{}{}{,}
dengan $\beta_n$ menyatakan
volume
bola satuan berdimensi $n$.

}
{

Kita menerapkan
Satz von Stokes
pada bola-bola
\mathl{B^n(P, \epsilon)}{} dengan batas
\mathl{S^{n-1}(P, \epsilon)}{}
\zusatzklammer {dengan pusat $P$ dan jari-jari $\epsilon$} {} {}
dan memperoleh

\mavergleichskettealign
{\vergleichskettealign
{ \operatorname{lim}_{   \epsilon \rightarrow  0  } \, { \frac{   1 }{  \epsilon^n } }     \int_{S^{n-1}(P, \epsilon)} \omega
}
{ =} { \operatorname{lim}_{   \epsilon \rightarrow  0  } \, { \frac{   1 }{  \epsilon^n } }     \int_{ \partial B^{n}(P, \epsilon)}   \omega
}
{ =} { \operatorname{lim}_{   \epsilon \rightarrow  0  } \, { \frac{   1 }{  \epsilon^n } }     \int_{B^{n}(P, \epsilon)} d \omega
}
{ =} { \operatorname{lim}_{   \epsilon \rightarrow  0  } \, { \frac{   1 }{  \epsilon^n } }     \int_{B^{n}(P, \epsilon)} f d \lambda^n
}
{ =} {\beta_n \cdot f(P)
}
}
{}
{}{,}
dengan kesamaan terakhir menggunakan kekontinuan $f$.

 }



\inputaufgabeklausurloesung
{12}

{

Buktikan teorema Stokes untuk manifold berbatas.

}
{

\teilbeweis {}{}{}
{Misalkan

\mathbed {U_i} {}
 {i \in I} {}
 {} {} {} {,}
suatu
\definitionsverweis {penutup terbuka}{}{}
dari $M$ dengan
\definitionsverweis {peta-peta berorientasi}{}{,}
dan misalkan

\mathbed {h_j} {}
 {j \in J} {}
 {} {} {} {,}
suatu
\definitionsverweis {partisi kesatuan yang terdiferensialkan secara kontinu dan tersubordinasi pada penutup tersebut}{}{,}
yang ada menurut
Satz 22.10 (Differentialgeometrie (Osnabrück 2023)). Untuk setiap

\mavergleichskette
{\vergleichskette
{ P
}
{ \in }{ M
}
{  }{
}
{    }{
}
{   }{
}
}
{}{}{}
terdapat suatu lingkungan terbuka

\mavergleichskette
{\vergleichskette
{ P
}
{ \in }{ V_P
}
{ \subseteq }{ M
}
{    }{
}
{   }{
}
}
{}{}{}
sedemikian sehingga

\mavergleichskette
{\vergleichskette
{ h_j \vert_{V_P }
}
{ = }{ 0
}
{ }{
}
{ }{
}
{ }{
}
}
{}{}{}
berlaku kecuali untuk berhingga banyak indeks. Misalkan $Y$ dukungan dari $\omega$. Penutup

\mavergleichskette
{\vergleichskette
{ Y
}
{ \subseteq }{ \bigcup_{P \in M} V_P
}
{  }{
}
{    }{
}
{   }{
}
}
{}{}{}
karena
\definitionsverweis {kekompakan}{}{}
yang diasumsikan, mempunyai suatu subpenutup berhingga, katakanlah

\mavergleichskettedisp
{\vergleichskette
{ Y
}
{ \subseteq} { V_{1} \cup \ldots \cup V_{r}
}
{ =} { V
}
{ } {
}
{ } {
}
}
{}{}{.}
Oleh karena itu, hanya berhingga banyak $h_j$ pada $V$ yang berbeda dari $0$. Kita menetapkan

\mavergleichskette
{\vergleichskette
{ \omega_j
}
{ = }{ h_j \omega
}
{  }{
}
{    }{
}
{   }{
}
}
{}{}{;}
bentuk-bentuk diferensial ini juga terdiferensialkan secara kontinu. Dukungan $h_j$ merupakan
\definitionsverweis {himpunan bagian tertutup}{}{}
dari $M$ yang terletak di dalam
\mathl{U_{i(j)}}{}. Karena itu, dukungan $\omega_j$ terletak di dalam
\mathl{Y \cap U_{i(j)}}{} dan kompak menurut
Aufgabe 13.10 (Differentialgeometrie (Osnabrück 2023)).
Berlaku

\mavergleichskettedisp
{\vergleichskette
{ \omega
}
{ =} { \sum_{j \in J} \omega_j
}
{ } {
}
{ } {
}
{ } {
}
}
{}{}{,}
dan hanya berhingga banyak bentuk diferensial
\mathl{\omega_j}{} yang berbeda dari $0$ karena

\mavergleichskette
{\vergleichskette
{ \omega_j \vert_{M \setminus Y}
}
{ = }{ 0
}
{ }{
}
{ }{
}
{ }{
}
}
{}{}{}
untuk setiap $j$, serta

\mavergleichskettedisp
{\vergleichskette
{ \omega_j \vert_V
}
{ =} { 0
}
{ } {
}
{ } {
}
{ } {
}
}
{}{}{}
untuk setiap $j$ kecuali berhingga banyak indeks. Berdasarkan
Additivität des Integrals von Differentialformen
dan
Additivität der äußeren Ableitung
pernyataan dapat dibuktikan secara terpisah untuk setiap $\omega_j$. Jadi, kita dapat menganggap bahwa diberikan suatu bentuk diferensial-$(n-1)$ yang terdiferensialkan secara kontinu dan mempunyai dukungan kompak yang seluruhnya terletak dalam suatu lingkungan peta $U$.}
{}
\teilbeweis {}{}{}
{Terdapat suatu
\definitionsverweis {difeomorfisme}{}{}
\maabb {\alpha} { U } { V
} {}
dengan

\mavergleichskette
{\vergleichskette
{ V
}
{ \subseteq }{ H
}
{ }{
}
{ }{
}
{ }{
}
}
{}{}{}
terbuka, yang sekaligus menginduksi suatu difeomorfisme di antara batas-batas

\mavergleichskette
{\vergleichskette
{ \partial U
}
{ = }{ \partial M \cap U
}
{  }{
}
{    }{
}
{   }{
}
}
{}{}{}
dan

\mavergleichskette
{\vergleichskette
{ \partial V
}
{ = }{ \partial H \cap V
}
{  }{
}
{    }{
}
{   }{
}
}
{}{}{}
. Berlaku

\mavergleichskettedisp
{\vergleichskette
{
\int_{ \partial U } \omega
}
{ = } {
\int_{ \partial V } \alpha^{-1 *}\omega
}
{ } {
}
{ } {
}
{ } {
}
}
{}{}{}
dan

\mavergleichskettedisp
{\vergleichskette
{
\int_{ U  }  d \omega
}
{ =} {
\int_{ V  }  \alpha^{-1 *} (  d \omega)
}
{ =} {
\int_{  V  }  d \left(  \alpha^{-1 *}\omega  \right)
}
{ } {
}
{ } {
}
}
{}{}{}
menurut
Lemma 20.2 (Differentialgeometrie (Osnabrück 2023))  (5).}
{\leerzeichen{}Jadi, kita dapat mulai dengan suatu bentuk diferensial yang terdiferensialkan secara kontinu, mempunyai dukungan kompak, dan didefinisikan pada

\mavergleichskette
{\vergleichskette
{ V
}
{ \subseteq }{ H
}
{ }{
}
{ }{
}
{ }{
}
}
{}{}{}
serta memperluasnya menjadi bentuk nol pada seluruh $H$ di luar dukungannya.}
\teilbeweis {}{}{}
{Karena kekompakan dukungan, terdapat suatu balok yang cukup besar

\mavergleichskette
{\vergleichskette
{ Q
}
{ \subseteq }{ H
}
{ }{
}
{ }{
}
{ }{
}
}
{}{}{,}
yang salah satu sisinya $S$ terletak pada $\partial H$ dan yang berpotongan dengan dukungan $\omega$ hanya di $S$. Pada semua sisi lain dari $Q$, bentuk $\omega$
\zusatzklammer {dan dengan demikian juga $d\omega$} {} {}
merupakan bentuk nol. Oleh karena itu, di satu pihak

\mavergleichskettedisp
{\vergleichskette
{
\int_{ H } d\omega
}
{ =} {
\int_{ Q } d\omega
}
{ } {
}
{ } {
}
{ } {
}
}
{}{}{}
dan di pihak lain

\mavergleichskettedisp
{\vergleichskette
{
\int_{ \partial H } \omega
}
{ = } {
\int_{ S } \omega
}
{ = } {
\int_{ \partial Q } \omega
}
{ } {
}
{ } {
}
}
{}{}{.}
Dengan demikian, pernyataan tersebut mengikuti dari Quaderversion des Satzes von Stokes.}
{}

 }
